This chapter pins down the concrete method developed in this
thesis: the predictor architecture
(Section~\ref{sec:method-architecture}), the training pipeline
that fits its parameters (Section~\ref{sec:method-training}),
the inference algorithm used at evaluation time
(Section~\ref{sec:method-inference}), the baselines
that the predictor is compared against
(Section~\ref{sec:method-baseline}), and the Python library
that implements all of the above
(Section~\ref{sec:method-implementation}).

\section{Overview}
\label{sec:method-overview}

The proposed predictor follows the neuro-symbolic architecture
of Section~\ref{sec:nn-features}: a neural backbone supplies
the per-cell unary potentials of a Markov Network, and a
shared pairwise matrix is learned jointly with the backbone
through a single structured loss.

Two design choices control the pipeline end-to-end:
\begin{itemize}
    \item \textbf{Graph structure} --- determined by the
    task. Chain graphs (the HMC dataset of
    Section~\ref{sec:datasets-hmc}) admit exact MAP inference
    via Viterbi (Section~\ref{subsec:viterbi}), so the
    structured-hinge loss with a Viterbi loss-augmented oracle
    can be used during training. Cyclic graphs (the Sudoku
    dataset of Section~\ref{sec:datasets-sudoku}) make MAP
    NP-hard; training switches to the LP-relaxed
    M\textsuperscript{3}N loss of
    Section~\ref{subsec:lp-m3n}, and inference at evaluation
    time falls back to an approximate decoder (Belief
    Propagation, presented in
    Section~\ref{sec:method-inference}).
    \item \textbf{Input modality} --- determined by the
    dataset. Symbolic inputs use a one-hot encoding and a
    linear backbone; visual inputs use a stack of MNIST
    images and the LeNet-5 backbone of
    Section~\ref{subsec:method-architecture-backbones}. The
    same trainer, optimizer, and inference path are reused
    across modalities; only the architecture YAML changes.
\end{itemize}

Two baselines are used for comparison
(Section~\ref{sec:method-baseline}): a two-stage OCR + solver
pipeline for visual Sudoku, which isolates the contribution of
joint structured training by replacing the M\textsuperscript{3}N
decision layer with a hard-coded solver while reusing the same
backbone; and a Forward-Backward marginal classifier on symbolic
HMC, which serves as the Bayes-optimal oracle floor in the chain
tightness experiments. The proposed M\textsuperscript{3}N
predictor is compared against both in
Chapter~\ref{chap:experiments}.

Every experiment in this thesis is driven by a YAML file
(Section~\ref{subsec:method-yaml}); the corresponding YAMLs
are listed in Chapter~\ref{chap:experiments}, and the
artifact format described in
Section~\ref{subsec:method-reproducibility} guarantees that
any reported number can be reproduced from a single
\texttt{config.yaml}.